\documentclass{article}

% Compile with: xelatex <file>.tex
\usepackage{fontspec}          % Unicode/system fonts (XeLaTeX/LuaLaTeX)
\usepackage[table]{xcolor}
\usepackage{amsmath, amssymb}  % Math
\newtheorem{theorem}{Theorem}
\usepackage{geometry}
\usepackage{graphicx}
\usepackage{titling}
\usepackage{fancyhdr}
\usepackage{lipsum}
\usepackage{xunicode}
\usepackage{longtable}
\usepackage{caption}
\usepackage{multirow}
\usepackage{xltxtra}% Enable line breaks for Thai text
\usepackage{tabularx,array} % in preamble
\usepackage{amsthm}
\XeTeXlinebreaklocale "th"
\XeTeXlinebreakskip = 0pt plus 2pt minus 1pt
\definecolor{graycell}{gray}{0.85}
\newtheorem{definition}{Definition}
\usepackage[style=ieee,sorting=nyt,maxbibnames=6,minbibnames=3,doi=false,isbn=false,url=false,eprint=false]{biblatex}
\addbibresource{refs.bib}
% Set up Thai fonts
\setmainfont[%
    ItalicFont={Laksaman-Italic.otf},%
    BoldFont={Laksaman-Bold.otf},%
    BoldItalicFont={Laksaman-BoldItalic.otf},%
    Script=Thai,%
    Scale=MatchLowercase,%
    WordSpace=1.25,%
    Mapping=tex-text,%
]{Laksaman.otf}
% Enable line breaks for Thai text
\XeTeXlinebreaklocale "th"
\XeTeXlinebreakskip = 0pt plus 2pt minus 1pt
\usepackage{setspace}

\geometry{
  a4paper,
  total={170mm,257mm},
  left=20mm,
  top=20mm,
}


% --- Title/author/date ---
\title{The Project Proposal of Course 2301399 Academic Year 2025}
\author{Peeradon Sarnkaew}
\date{September 2025}

% --- Header/footer ---
\pagestyle{fancy}
\fancyhf{}
\fancyhead[L]{The Project Proposal of Course 2301399 Academic Year 2025}
\fancyhead[R]{\theauthor}
\fancyfoot[L]{\today}
\fancyfoot[R]{\includegraphics[width=2cm]{cu.png}}
\renewcommand{\headrulewidth}{0.4pt}
\renewcommand{\footrulewidth}{0pt}

% (Optional) custom \maketitle without date
\makeatletter
\def\@maketitle{%
  \newpage
  \null
  \vskip 1em%
  \begin{center}%
    {\LARGE \@title \par}%
    \vskip 1em%
  \end{center}%
  \par
  \vskip 1em}
\makeatother

\begin{document}

\maketitle

\noindent
\begin{tabularx}{\textwidth}{@{}l>{\raggedright\arraybackslash}X@{}}
  Project Title (Thai)   & การเรียนรู้การเข้ารหัสเวลาสัมบูรณ์และสัมพัทธ์ สำหรับกราฟพลวัตเชิงเวลาแบบต่อเนื่อง \\
  Project Title (English) & KAN-MAMMOTE: Learnable Absolute and Relative Time Encoding for Continuous-Time Dynamic temporal Graphs \\
  By & \theauthor \\
     & Computer Science Program,\\
     & Department of Mathematics and Computer Science, Faculty of Science, \\
     & Chulalongkorn University \\
  Advisor & Asst. Prof. Naruemon Pratanwanich \\
  Co-Advisor & Senior Lecturer Natthawut Kertkeidkachorn \\
  & Chief Senior Researcher XIN, Liu \\
  \\
\end{tabularx}

\hrule
\section*{Background and Rationale}

\paragraph{Motivation and Background}
Interactions within complex systems, from social networks to financial markets, are inherently dynamic and occur in continuous time. These systems are naturally modeled as Continuous-Time Dynamic Graphs (CTDGs), where the graph topology evolves asynchronously over a real-valued timeline. To faithfully capture the underlying dynamics, machine learning frameworks must jointly model the graph structure and the precise temporal information embedded in each event. The irregular sampling of these events presents a fundamental challenge; without a robust method for encoding time, models fail to capture essential temporal patterns such as activity bursts, seasonal cycles, and memory decay effects, which are critical for tasks like temporal link prediction, event forecasting, and time-aware recommendation.

\paragraph{Limitations of Prior Work}
To address this temporal component, significant research has focused on Functional Time Encoding (FTE). Early methods such as Time2Vec \cite{kazemi2019time2vec} and kernel-based approaches \cite{xu2019functional} transform raw timestamps into vector representations. While effective in certain contexts, these methods are constrained in their expressive power. Time2Vec, for instance, relies on a fixed combination of linear and sinusoidal functions, rendering it insufficient for modeling complex, non-periodic temporal behaviors. Similarly, Mercer and Bochner-based encoders are often restricted to capturing periodic phenomena through sine and cosine bases. More recent architectures like LeTE \cite{lete_placeholder} have advanced the field by employing a mixture of spline and Fourier modules to model locality and periodicity in absolute time. However, despite these improvements, fundamental limitations persist, preventing a complete and unified understanding of temporal dynamics.

\paragraph{The Duality of Time: A Core Unresolved Problem}
A critical gap in the existing literature is the incomplete treatment of time's dual nature: its \textit{absolute} and \textit{relative} aspects. The majority of contemporary methods focus on only one of these facets. Models that process relative time via time gaps ($\Delta t$) excel at capturing local, translation-invariant dynamics like memory decay. However, they are fundamentally incapable of modeling global, context-dependent patterns. For example, such models cannot distinguish the significance of a one-day gap between December 31st and January 1st from one between two days in mid-July, thereby losing crucial information tied to calendar-based seasonality. Conversely, models focusing solely on absolute timestamps can capture global regimes but struggle to precisely model the influence of event recency and inter-event duration. Furthermore, existing encoders are typically constrained by a single family of basis functions (e.g., sinusoidal, spline), which restricts their capacity to approximate the diverse array of temporal patterns present in real-world data. A truly robust framework must therefore not only represent both absolute and relative time but also possess the functional flexibility to capture their distinct characteristics and, most importantly, model their intricate interactions.

\paragraph{Our Proposal: KAN-MAMMOTE}
To overcome these challenges, we propose \textbf{KAN-MAMMOTE}, a dual-stream, learnable framework designed to create a unified representation of absolute and relative time. Our approach is built on two novel, specialized components. For absolute time, we introduce \textit{K-MOTE} (Kolmogorov-Arnold Mixture-of-Time-Experts), a flexible module that leverages a diverse set of Kolmogorov-Arnold Networks (KANs)—based on Fourier, B-spline, and Wavelet functions—to capture a wide spectrum of global temporal patterns. For relative time, we employ a learnable \textit{SM-Kernel}, a Spectral Mixture kernel from Gaussian Process theory that is provably translation-invariant and adept at modeling local dynamics such as temporal decay and oscillations. The core innovation of our work lies in the fusion of these two streams. We introduce a modified Mamba architecture, a type of State Space Model (SSM), where the absolute and relative time embeddings are used to dynamically generate its selective state-update parameter, $\Delta$. This transforms Mamba's selection mechanism into a sophisticated temporal gate, allowing the model to reason about how an event's influence decays over time based on the absolute temporal context. This fused representation, which we term an absolute-relative time embedding, is then injected into downstream GNN backbones to replace its original Functional Time Encoder Model. Our design not only generalizes prior encoding strategies but also provides mathematical guarantees of key properties, enabling robust and interpretable temporal learning in CTDGs.

\paragraph{Contributions and Paper Structure}
In this paper, we systematically develop and validate a novel framework for temporal representation learning in CTDGs. Our primary contributions are threefold: (1) We propose \textbf{KAN-MAMMOTE}, a novel dual-stream architecture that, for the first time, jointly models absolute and relative time using highly flexible and mathematically principled function approximators. (2) We introduce \textit{K-MOTE}, a mixture-of-experts module based on Kolmogorov-Arnold Networks, for expressive absolute time encoding. (3) We design a novel fusion mechanism that repurposes a State Space Model (Mamba) into a dynamic temporal gate, allowing the model to learn the complex interplay between absolute and relative temporal signals. The remainder of this paper is structured to rigorously detail and validate these contributions. We begin by situating our work within the broader context of temporal graph learning before formally presenting the KAN-MAMMOTE architecture and its theoretical underpinnings. The core of our validation lies in an extensive empirical evaluation, demonstrating that KAN-MAMMOTE significantly outperforms state-of-the-art baselines on several real-world CTDG benchmarks. To dissect the sources of this performance gain, we conduct a comprehensive suite of ablation studies that quantify the individual contributions of the dual-stream design, the mixture-of-experts module, and our dynamic fusion mechanism. Crucially, we address the challenge of model interpretability through a qualitative analysis on a synthetic dataset, where we visualize the learned functions of K-MOTE and the SM-Kernel to demonstrate their ability to recover known ground-truth temporal patterns, such as seasonality and memory decay. Furthermore, we analyze the model's computational efficiency, providing both theoretical complexity and empirical runtime benchmarks. Finally, we conclude by discussing the broader implications and limitations of our framework and outline promising directions for future research in dynamic representation learning.
\section*{Objective}
The primary objective of this work is to design, implement, and validate KAN-MAMMOTE, a novel dual-stream temporal encoding framework for CTDGs
\section*{Detailed Objectives}
To develop a dual-stream, learnable time encoding framework for CTDGs that captures the distinct and interactive nature of absolute and relative temporal signals, leading to superior performance and interpretability. The core objectives are to:
\begin{itemize}
  \item \textbf{Develop K-MOTE:} A novel mixture-of-time-experts module for absolute time encoding, leveraging a diverse set of Kolmogorov-Arnold Network (KAN) bases (Fourier, B-spline, Wavelet) to flexibly capture complex global patterns like seasonality and trends.
  \item \textbf{Implement the SM-Kernel:} A stationary Spectral Mixture kernel for relative time encoding that is provably translation-invariant, designed to accurately model local temporal dynamics such as memory decay and event burstiness.
  \item \textbf{Design a Dynamic Temporal Fusion Mechanism:} A novel fusion strategy that repurposes a Mamba (State Space Model) architecture. The mechanism will use the learned absolute and relative time embeddings to dynamically control Mamba's internal selection parameter ($\Delta'$), effectively creating a time-aware gating system.
  \item \textbf{Conduct Rigorous Empirical Validation:} Systematically evaluate KAN-MAMMOTE against state-of-the-art baselines on a comprehensive suite of real-world CTDG benchmarks, and perform extensive ablation studies to verify the contribution of each architectural component.
\end{itemize}

\section*{Scope}
\begin{itemize}
  \item \textbf{Problem Setting:} The project focuses on the CTDG setting where node interactions are asynchronous and timestamped with real-valued time, targeting the downstream task of dynamic link prediction.
  \item \textbf{Model-Agnostic Design:} The KAN-MAMMOTE module is designed to be model-agnostic, functioning as a plug-and-play replacement for the native time encoding components (e.g., `TimeEncoder` modules or direct `$\Delta$t` inputs) in established GNN backbones, including TGAT, TGN, TCL, DyGFormer, and DyGMamba.
  \item \textbf{Evaluation Protocol:} Evaluations will be conducted using the rigorous protocol and datasets from the "Towards Better Evaluation for Dynamic Link Prediction" benchmark (Poursafaei et al., 2022), which includes Wikipedia, MOOC, LastFM, Reddit, SocialEvo, and UCI. Performance will be measured using standard metrics for the task: Area Under the ROC Curve (AUC) and Average Precision (AP).
  \item \textbf{Training:} The time encoding framework will be trained end-to-end jointly with the GNN backbone and the downstream link prediction objective.
\end{itemize}

\section*{Project Activities}
\begin{enumerate}
  \item \textbf{Literature Review:} Conduct a comprehensive survey of existing time encoding methods for dynamic graphs and related domains, including FTR, LeTE, CTLPE, and Time2Vec.
  \item \textbf{Component Implementation:}
    \begin{itemize}
        \item Implement the \textbf{K-MOTE} module using multiple KAN function approximators.
        \item Implement the \textbf{SM-Kernel} module for relative time, leveraging scalable Gaussian Process libraries.
    \end{itemize}
  \item \textbf{Fusion Module Development:} Implement the dynamic fusion module by modifying a Mamba (SSM) architecture to accept the absolute and relative embeddings as controllers for its selection mechanism.
  \item \textbf{Integration and Baseline Comparison:} Integrate the full KAN-MAMMOTE encoder into selected CTDG models (e.g., TGAT, TGN) and run comprehensive baseline comparisons on the specified benchmarks.
  \item \textbf{Ablation Studies:} Conduct extensive ablation experiments to isolate the impact of: (a) the dual-stream design vs. single-stream variants, (b) the choice of K-MOTE experts, and (c) our dynamic Mamba-based fusion vs. simpler fusion methods.
  \item \textbf{Qualitative and Quantitative Analysis:}
    \begin{itemize}
        \item Analyze the interpretability of the model on a synthetic dataset by visualizing the learned functions of the K-MOTE experts and the learned power spectrum of the SM-Kernel.
        \item Analyze the model's computational complexity, empirical runtime, and generalization capabilities.
    \end{itemize}
  \item \textbf{Documentation and Dissemination:} Document the framework, prepare a manuscript for submission to a top-tier conference, and release the code to ensure reproducibility.
\end{enumerate}

\section*{Timeline of the study}

\renewcommand{\arraystretch}{1.25}
\setlength{\tabcolsep}{6pt}        % was 8pt; tighter padding
\emergencystretch=2em               % relax line-breaking to avoid overfull boxes
\hyphenation{SpatioTemporal KAN-MAMMOTE}

\begin{tabular}{| >{\raggedright\arraybackslash}p{6.8cm} | *{7}{>{\centering\arraybackslash}p{0.95cm}|}}
\hline
\textbf{Activities} & \textbf{May} & \textbf{Jun} & \textbf{Jul} & \textbf{Aug} & \textbf{Sep} & \textbf{Oct} & \textbf{Nov} \\
\hline
1. Research \& Study The Foundation of Graph Neural Networks.
    & \cellcolor{graycell} &  &  &  &  &  &   \\
\hline
2. Research Continuous-Time Dynamic SpatioTemporal Graph Neural Networks
    &  & \cellcolor{graycell} &  &  &  &  &   \\
\hline
3. Research existing Time Encoding Models, Sequence Analysis Theory, Related Sequence Analysis Methods and Math Functions
    &  & \cellcolor{graycell} & \cellcolor{graycell} &  &  &  &   \\
\hline
4. Develop baseline model for link prediction task, and develop K-MOTE structure
    &  &  & \cellcolor{graycell} &  &  &  &   \\
\hline
5. Develop SM-Kernel structure
    &  &  &  & \cellcolor{graycell} &  &  &   \\
\hline
6. Develop fusion strategy with Mamba layer and integrate all components into KAN-MAMMOTE
    &  &  &  & \cellcolor{graycell} & \cellcolor{graycell} &  &   \\
\hline
7. Evaluate against baseline models
    &  &  &  &  & \cellcolor{graycell} & \cellcolor{graycell} &  \\
\hline
8. Conduct experiments to ensure interpretability and mathematically prove the required theory for each component
    &  &  &  &  &  & \cellcolor{graycell} &  \\
\hline
9. Conduct ablation study and time complexity analysis
    &  &  &  &  &  & \cellcolor{graycell} & \cellcolor{graycell} \\
\hline
10. Write and review project report
    &  &  &  &  &  & \cellcolor{graycell} & \cellcolor{graycell} \\
\hline
\end{tabular}



\section*{Benefits}
\subsection*{A. Benefit for the Implementor}
The implementor will gain significant expertise at the intersection of temporal graph modeling, modern deep learning architectures (KANs, SSMs), and kernel methods. This project will strengthen skills in architectural design, rigorous experimental validation, and preparing research for publication in the machine learning community.

\subsection*{B. Benefit for Users and the Research Community}
Users of temporal GNN systems will benefit from a powerful, plug-and-play module that can substantially improve predictive accuracy and robustness on complex event data. For the research community, this work will introduce a novel, interpretable, and mathematically principled framework for temporal encoding, potentially setting a new standard for modeling time in dynamic systems.

\section*{Equipment}
\begin{itemize}
    \item \textbf{Primary Computing:} Access to a GPU-equipped server (HPC Kagayaki) with a standard Python environment (PyTorch, PyG, GPyTorch).
    \item \textbf{Local Development Machine:} Computer with CPU Intel(R) Core(TM) i5-13500HX, 16GB RAM, and GPU NVIDIA GeForce RTX4060 (8GB VRAM).
    \item \textbf{Cloud Computing (Supplemental):} Google Colaboratory Pro+ subscription for supplemental experimentation and prototyping.
\end{itemize}

\section*{Budget}
\begin{itemize}
    \item Google Colaboratory Service (Pro+ Subscription, 1 month): \~1,700 THB
    \item Peripheral Hardware (LG 24" QHD Monitor, SSD, Keyboard): \~7,400 THB
    \item \textbf{Total Estimated Budget:} \~9,100 THB
\end{itemize}

\section*{References}
\textit{(Note: This section should be expanded to include seminal papers for the methods you are building upon.)}
\begin{itemize}
    \item \textbf{[Your References]:}
    \item Kim et al., "Continuous-Time Linear Positional Embedding for Irregular Time Series Forecasting," 2024.
    \item Ding et al., "Rethinking Time Encoding via Learnable Transformation Functions," ICLR 2024.
    \item Huang et al., "Self-Attention with Functional Time Representation Learning," 2023.
    \item \textbf{[Suggested Additions]:}
    \item Gu, A., \& Dao, T. (2023). Mamba: Linear-Time Sequence Modeling with Selective State Spaces. \textit{arXiv preprint arXiv:2312.00752}.
    \item Liu, Z., et al. (2024). KAN: Kolmogorov-Arnold Networks. \textit{arXiv preprint arXiv:2404.19756}.
    \item Wilson, A. G., \& Adams, R. P. (2013). Gaussian process kernels for pattern discovery and extrapolation. In \textit{International conference on machine learning}.
    \item Rossi, E., et al. (2020). Temporal Graph Networks for Deep Learning on Dynamic Graphs. \textit{arXiv preprint arXiv:2006.10637}.
    \item Poursafaei, F., et al. (2022). Towards Better Evaluation for Dynamic Link Prediction on Graphs. \textit{Advances in Neural Information Processing Systems (NeurIPS)}.
\end{itemize}
\printbibliography[title={References}]


\newpage

\section{Problem Formulation}

In this section, we formally define a Continuous-Time Dynamic Spatiotemporal Graph (CTDSG), the task of link prediction within this graphical structure, and the incorporation of temporal embeddings as augmented features for nodes and edges.

\subsection{Continuous-Time Dynamic Spatiotemporal Graph}

A Continuous-Time Dynamic Spatiotemporal Graph is a graph where the topological structure and the attributes of its components evolve continuously over time and are associated with spatial properties. Unlike discrete-time dynamic graphs that are represented as a sequence of static snapshots, a continuous-time graph is defined by a stream of events, each with a precise timestamp. [7, 19, 23]

\begin{definition}[Continuous-Time Dynamic Spatiotemporal Graph]
A \textbf{Continuous-Time Dynamic Spatiotemporal Graph} (CTDSG) is a tuple $\mathcal{G} = (\mathcal{V}, \mathcal{E}, \mathcal{S}, \mathcal{T})$, where:
\begin{itemize}
    \item $\mathcal{V}$ is the set of vertices (or nodes).
    \item $\mathcal{E} \subseteq \mathcal{V} \times \mathcal{V} \times \mathbb{R}^+ \times \mathcal{X}_e$ is a multiset of timed and attributed edges. An edge $e = (u, v, t, \mathbf{x}_e) \in \mathcal{E}$ represents an interaction between node $u$ and node $v$ that occurs at a continuous time $t \in \mathcal{T}$, with $\mathbf{x}_e \in \mathcal{X}_e$ being a vector of edge features.
    \item $\mathcal{S}: \mathcal{V} \to \mathbb{R}^d$ is a function that maps each node to a $d$-dimensional spatial coordinate or embedding. These spatial attributes can be static or dynamic. In the dynamic case, $\mathcal{S}$ can be represented as a function of time, $\mathcal{S}(v, t)$.
    \item $\mathcal{T}$ is the continuous time domain, typically a subset of $\mathbb{R}^+$.
\end{itemize}
The set of nodes $\mathcal{V}$ can also be dynamic, with nodes being added or removed at specific timestamps. Similarly, nodes can have time-varying features, denoted by a function $\mathbf{x}_v(t) \in \mathcal{X}_v$.
\end{definition}

\subsection{Link Prediction in a Continuous-Time Dynamic Spatiotemporal Graph}

The primary objective of link prediction in a CTDSG is to forecast the emergence of a future edge between a pair of nodes, given the history of the graph's evolution up to a certain point in time. [9, 13] This task is crucial for understanding and modeling the dynamics of evolving networks.

\begin{definition}[Link Prediction Task]
Given a Continuous-Time Dynamic Spatiotemporal Graph $\mathcal{G}$ observed over the time interval $[0, T)$, the task of link prediction is to learn a function $f$ that predicts the probability of a new edge $(u, v)$ forming at a future time $t' > T$:

$$ f: (\mathcal{V}, \mathcal{E}_{[0, T)}, \mathcal{S}, u, v, t') \mapsto p(e_{u,v,t'} \in \mathcal{E}) $$

where $\mathcal{E}_{[0, T)}$ represents the set of all edges that occurred in the time interval $[0, T)$, and $p(e_{u,v,t'} \in \mathcal{E})$ is the probability of the existence of an edge between nodes $u, v \in \mathcal{V}$ at time $t'$. This is often formulated as a binary classification problem where the goal is to distinguish between future true edges and non-existent edges.
\end{definition}
\subsection{The Use of Time Encoding in CTDSGs}

The timestamps associated with events in a Continuous-Time Dynamic Graph are raw scalar values (e.g., UNIX timestamps). To be utilized by deep learning models, this temporal information must be transformed into a high-dimensional vector representation through a \textit{time encoding} function, denoted as $\mathcal{TE}$. This process generates a time embedding that captures the underlying temporal characteristics in a format amenable to neural network processing.

In the domain of dynamic graph representation learning, the standard practice is to concatenate this time embedding with existing node and/or edge features. This creates an augmented feature vector that serves as the input to the graph neural network backbone for a given interaction $e_k = (u, v, t_k, \mathbf{x}_{e_k})$. This allows for flexibility in the dimensionality of the time embedding. The resulting input vector $\mathbf{x}_k$ is formally expressed as:
\begin{equation}
\mathbf{x}_k = \text{NodeFeatures}(u, v, t_k) \Vert \mathbf{x}_{e_k} \Vert \mathcal{TE}(\cdot) \in \mathbb{R}^{d_n + d_e + d_t}
\end{equation}
where $\Vert$ denotes the concatenation operation, and $d_n$, $d_e$, and $d_t$ represent the dimensions of the node features, edge features, and the time embedding, respectively.

A critical design choice lies in the input to the time encoding function $\mathcal{TE}(\cdot)$. Many contemporary CTDG models, including influential baselines like TGAT and TGN, simplify this by using only the \textit{relative time difference} ($\Delta t$) between the current event and a previous event. While this approach effectively models local, translation-invariant dynamics, it inherently discards the absolute temporal context. As detailed in our rationale, this limitation prevents such models from capturing global patterns like seasonality or time-of-day effects.

To overcome this fundamental trade-off, our work redefines the nature of the temporal embedding. The \textbf{KAN-MAMMOTE} framework is designed to produce a single, unified temporal representation, $\mathbf{T}_k$, that is a function of \textit{both} the absolute timestamp of the current event, $t_k$, and the relative time difference to the previous event, $\Delta t_k$. This fused embedding encapsulates the rich information from both temporal aspects. The final input to the GNN backbone is therefore constructed as:
\begin{equation}
\mathbf{x}_k = \text{NodeFeatures}(u, v, t_k) \Vert \mathbf{x}_{e_k} \Vert \mathbf{T}_k \quad \text{where} \quad \mathbf{T}_k = \text{KAN-MAMMOTE}(t_k, \Delta t_k)
\end{equation}
By defining the time encoding in this manner, we enable the downstream GNN to leverage a far more comprehensive temporal signal, allowing it to reason about events based on both their global context and their local recency.
\section{The KAN-MAMMOTE Framework}

To address the limitations of prior work, we introduce \textbf{KAN-MAMMOTE}, a novel time encoding framework designed to create a unified and expressive representation from both absolute and relative temporal signals. The architecture is founded on three core principles: (1) specialized, dual-stream encoding to capture the distinct nature of absolute and relative time; (2) the use of highly flexible and mathematically principled function approximators for each stream; and (3) a dynamic fusion mechanism that models the interaction between the two temporal aspects. An overview of the full architecture is presented in Figure~\ref{fig:kan-mammote-architecture}.

\begin{figure}[h!]
    \centering
    % You would include your KAN-MAMMOTE diagram here
    % \includegraphics[width=\textwidth]{path/to/kan-mammote-diagram.png}
    \caption{The KAN-MAMMOTE architecture. The model processes absolute time ($t_k$) and relative time ($\Delta t$) through two specialized streams, K-MOTE and SM-Kernel, respectively. The resulting embeddings ($\mathbf{u}_k, \mathbf{v}_k$) are used to dynamically modulate a Mamba (SSM) block, which processes the absolute time sequence to produce the final, unified temporal representation $\mathbf{T}_k$.}
    \label{fig:kan-mammote-architecture}
\end{figure}

\subsection{Dual-Stream Architecture for Temporal Encoding}
Our framework begins by explicitly separating the encoding of absolute and relative time into two parallel streams. For each event $k$ occurring at timestamp $t_k$, we compute the relative time difference $\Delta t_k = t_k - t_{k-1}$ with respect to the immediately preceding event. These two signals, $t_k$ and $\Delta t_k$, are then fed into their respective specialized encoders.

\subsection{Absolute Time Encoding: K-MOTE}
To capture the complex, non-stationary patterns inherent in absolute timestamps (e.g., seasonality, trends, time-of-day effects), we propose the \textbf{Kolmogorov-Arnold Mixture-of-Time-Experts (K-MOTE)} module. This module is built upon Kolmogorov-Arnold Networks (KANs), which feature learnable activation functions on their edges, offering superior flexibility and interpretability. K-MOTE employs a mixture-of-experts design, with specialized experts (Fourier-KAN, Spline-KAN, Wavelet-KAN) that are adaptively weighted by a gating network to capture a wide spectrum of temporal patterns. The module processes the sequence of absolute timestamps $t_k$ to produce a rich embedding $\mathbf{u}_k \in \mathbb{R}^D$:
\begin{equation}
    \mathbf{u}_k = \text{K-MOTE}(t_k)
\end{equation}

\subsection{Relative Time Encoding: SM-Kernel}
To model the relative time gap $\Delta t_k$, we employ a learnable \textbf{Spectral Mixture (SM) Kernel}. Derived from Gaussian Process theory, the SM-Kernel is both \textbf{provably translation-invariant}—correctly modeling the stationary nature of processes like memory decay—and highly interpretable, as it automatically discovers latent periodicities by learning a power spectral density. The SM-Kernel processes the time gap $\Delta t_k$ to produce the relative time embedding $\mathbf{v}_k \in \mathbb{R}^D$:
\begin{equation}
    \mathbf{v}_k = \text{SM-Kernel}(\Delta t_k)
\end{equation}

\subsection{Dynamic Temporal Fusion via Modulated Mamba}
The core innovation of our framework lies in how the absolute and relative embeddings are fused. Instead of a simple post-hoc combination, we propose a mechanism that dynamically \textbf{modulates} the internal state transitions of a Mamba2 architecture.

Mamba's primary strength is its selective State Space Model (SSM), which uses an input-dependent gate, $\Delta$, to control how its recurrent state persists or resets. In Mamba2, this gate, which we term $\mathbf{dt}_{\text{content}}$, is computed from the input sequence—in our case, the absolute time embeddings $\mathbf{u}_k$. This provides the model with content-awareness. Our strategy is to enhance, not replace, this mechanism by injecting our dual-stream temporal knowledge. This is achieved in three steps:

\begin{enumerate}
    \item \textbf{Content-Gate Calculation:} The sequence of absolute embeddings, $\mathbf{u}_k$, is passed to the Mamba2 block, which computes its internal, content-based gate $\mathbf{dt}_{\text{content}}$ for each event in the sequence.
    
    \item \textbf{Temporal-Gate Calculation:} In parallel, both the absolute embedding $\mathbf{u}_k$ and the relative embedding $\mathbf{v}_k$ are concatenated and passed through a dedicated \textbf{Fusion MLP}. This MLP outputs a multiplicative \textit{temporal gate}, $\mathbf{g}_t$, which is scaled by a factor of two and passed through a sigmoid function to be centered around 1.0.
    \begin{equation}
        \mathbf{g}_t = 2 \cdot \sigma\left(\text{MLP}([\mathbf{u}_k \Vert \mathbf{v}_k])\right)
    \end{equation}
    
    \item \textbf{Multiplicative Modulation:} The final, fused gate, $\mathbf{dt}_{\text{fused}}$, is computed by the element-wise product of Mamba's content-based gate and our temporal gate. This fused gate is then used to control the state update of the Mamba SSM at every step in the sequence.
    \begin{equation}
        \mathbf{dt}_{\text{fused}} = \mathbf{dt}_{\text{content}} \odot \mathbf{g}_t
    \end{equation}
\end{enumerate}

This modulation transforms the Mamba2 block into a time-aware sequence model. The learned temporal dynamics act as a controller for its memory, allowing it to amplify, suppress, or maintain its state based on a synergistic combination of sequence content and dual-stream temporal context. The final, unified temporal representation $\mathbf{T}_k \in \mathbb{R}^D$ is the output of this dynamically modulated process.
\end{document}